% ============================================================================
% markov_models_neutral.tex
%
% Catalogue of every Markov model specified in a doctoral thesis on
% stochastic branching and release processes. Neutral-terminology copy,
% prepared for external automated review: biological framing has been
% removed and the mathematical specifications retained in full.
%
% Compiled standalone:  pdflatex markov_models_neutral.tex  (run twice)
% ============================================================================

\documentclass[11pt,a4paper]{report}
\usepackage[a4paper,margin=1in]{geometry}
\usepackage[T1]{fontenc}
\usepackage[utf8]{inputenc}
\usepackage{lmodern}
\usepackage{microtype}
\usepackage{amsmath,amssymb,amsthm,mathtools}
\usepackage{booktabs}
\usepackage{array}
\usepackage{longtable}
\usepackage{enumitem}
\usepackage{xcolor}
\usepackage{hyperref}

\hypersetup{
  colorlinks=true,
  linkcolor=blue!55!black,
  urlcolor=blue!60!black,
  pdftitle={Catalogue of Markov model specifications}
}

% --- House notation ----------------------------------------------------------
\newcommand{\pr}[1]{\mathbb{P}\!\left(#1\right)}
\newcommand{\ex}[1]{\mathbb{E}\!\left(#1\right)}
\newcommand{\mean}[1]{\mathbb{E}\!\left[#1\right]}
\newcommand{\var}[1]{\operatorname{Var}\!\left(#1\right)}
\newcommand{\E}{\mathbb{E}}
\newcommand{\Prob}{\mathbb{P}}
\newcommand{\dd}{\mathrm{d}}
\newcommand{\ee}{\mathrm{e}}
\newcommand{\ddt}[1]{\frac{\dd #1}{\dd t}}
\newcommand{\dda}[2]{\frac{\dd #1}{\dd #2}}
\newcommand{\pddt}[1]{\frac{\partial #1}{\partial t}}
\newcommand{\pdd}[2]{\frac{\partial #1}{\partial #2}}
\newcommand{\lrb}[1]{\left(#1\right)}
\newcommand{\lrbs}[1]{\left[#1\right]}
\newcommand{\ind}{\mathbf 1}
\newcommand{\R}{\mathsf{R}}
\newcommand{\Ifix}{I_{\mathrm{fix}}}
\newcommand{\Kr}{\mathcal{K}}      % release size at rupture
\newcommand{\Act}{\mathcal{I}}     % active compartments
\newcommand{\Free}{\mathcal{V}}    % free particles
\newcommand{\Xt}{X_t}
\newcommand{\Wt}{W_t}
\newcommand{\Bint}{\mathcal{B}}
\newcommand{\Lap}[1]{\widetilde{#1}}

\newcommand{\source}[1]{\par\noindent{\small\textit{Specified in #1.}}\par\medskip}

\setcounter{secnumdepth}{2}
\setcounter{tocdepth}{1}
\allowdisplaybreaks

\begin{document}

\begin{titlepage}
\centering
\vspace*{4.2cm}
{\LARGE\bfseries Catalogue of Markov model specifications\par}
\vspace{10pt}
{\large Extracted from a doctoral thesis on stochastic branching and release
processes\par}
\vspace{1.6cm}
{\large Neutral-terminology copy for automated review\par}
\vspace{6pt}
{September 2026\par}
\vfill
{\small Every entry gives the thesis location, state space, transition rates,
parameters and attached conventions. Rates are transcribed verbatim from the
chapter sources.\par}
\vspace*{1.2cm}
\end{titlepage}

\tableofcontents
\clearpage

% ----------------------------------------------------------------------------
\chapter*{Scope and conventions}
\addcontentsline{toc}{chapter}{Scope and conventions}
% ----------------------------------------------------------------------------

This document lists every stochastic model in the source thesis whose law is
\emph{specified}: a state space plus a transition mechanism (transition
matrix, generator, jump rates, or transition kernel). Deterministic
population equations are included only where the thesis derives them as
exact projections of a stochastic model, and the entry says so. Derived
results are quoted only where they are part of the model's working
apparatus (characteristic roots, regime boundaries, working parameter
points).

Time is continuous and chains are time homogeneous unless an entry says
otherwise; $o(\Delta t)$ terms are suppressed. Per-capita rates multiply the
relevant count to give total rates. Model numbers M1--M31 are used for
cross-reference. Sections are cited by thesis chapter and section number
only.

% ----------------------------------------------------------------------------
\chapter{Catalogue at a glance}
% ----------------------------------------------------------------------------

{\small
\begin{longtable}{@{}p{0.055\textwidth}p{0.40\textwidth}p{0.32\textwidth}p{0.11\textwidth}@{}}
\toprule
& Model & State space & Where \\
\midrule
\endhead
M1 & Simple random walk & $\mathbb{Z}$ (discrete time) & \S2.2.1.1 \\
M2 & Galton--Watson process (general; binary case) & $\{0,1,2,\ldots\}$ (discrete time) & \S2.2.1.2; Ch.~3 \\
M3 & Poisson process & $\{0,1,2,\ldots\}$ & \S2.2.2.2 \\
M4 & Yule process & $\{1,2,\ldots\}$ & \S2.2.2.3 \\
M5 & Linear birth--death process (incl.\ time-inhomogeneous rates) & $\{0,1,2,\ldots\}$ & \S2.2.2.4; \S2.2.3 \\
M6 & Birth--death--catastrophe process & $\{0,1,2,\ldots\}\cup\{H\}$ & \S2.2.2.5; \S4.2; \S5.2.2 \\
M7 & Seasonal two-population process & $\{0,1,\ldots\}^2$ & \S2.2.3.1 \\
M8 & Logistic-environment type-count process & $\{1,2,\ldots\}$, driven by an ODE & \S2.2.3.2; \S9.2 \\
M9 & Gated release (coupled ODE--CTMC) & $[0,\infty)^2\times\{0,1\}$ & \S2.2.4.1 \\
M10 & Surging gated release & $[0,\infty)^2\times\{0,1\}$ & \S2.2.4.2 \\
M11 & Absorption-only model & $\{0,\ldots,N\}$ pairs & \S2.B.1 \\
M12 & Absorption--death model & compartment pairs & \S2.3.3.2 \\
M13 & Absorption--birth--death model & compartment pairs & \S2.B.2 \\
M14 & Chained immediate transfer ($\mu=0$ birth--catastrophe chain) & one load per compartment, $M$ compartments & \S5.10; \S8.2.1 \\
M15 & Budding comparator compartment & two independent event streams & \S5.11; \S6.2.2 \\
M16 & Standard two-variable population model (deterministic projection) & $(\Act,\Free)$ ODEs & \S1.5; \S6.2.1 \\
M17 & Kernel-driven renewal system & age-structured means & \S6.3--6.4 \\
M18 & Branching approximation of the early phase; flooding parameter $L$ & offspring counts & \S6.6 \\
M19 & Reset process (recurrent release) & $\{0,1,2,\ldots\}$ + compartment removal & \S6.7.1 \\
M20 & Partial release (boolean thinning) & mean-field stock system & \S6.7.2 \\
M21 & Budding-kernel extension with sigmoid incidence & kernel pair + incidence & \S6.7.3 \\
M22 & Two-type birth--death--conversion with type-specific catastrophe & $\mathbb{Z}_{\ge0}^2\cup\{(\R,\R)\}$ & \S7.2 \\
M23 & Replicator--suppressor process & $\{0,\ldots\}\times\{0,\ldots,N\}$ & \S8.3.2 \\
M24 & Fixed-removal extension of the chained process & release sizes + budget & \S8.4 \\
M25 & Competitive-reversal chain (evolving resistance) & $\{0,\ldots,N\}$ + resistances & \S9.3.2 \\
M26 & Fitness-coupled type-origination variant & populations + fitness + type counts & \S9.3.6 \\
M27 & Rise--run--ruin--rejuvenation (RRRR) & count + potential $(\Xt,V(t))$ & \S9.4.2 \\
M28 & Multiplicative-dissipation automaton & $L\times L$ lattice + potentials & \S9.5 \\
M29 & Birth--conversion model (incl.\ vitality variant) & $\{0,1,\ldots\}^2$ & \S9.A \\
M30 & Public-good accumulation process & count + good concentration & \S9.B \\
M31 & Three-environment birth--death ordering & three birth--death processes & \S10.3.1 \\
\bottomrule
\end{longtable}
}

% ----------------------------------------------------------------------------
\chapter{Foundational processes (thesis Chapter 2)}
% ----------------------------------------------------------------------------

\section{M1: Simple random walk}
\source{\S2.2.1.1}

\paragraph{State space.} $\mathbb{Z}$, discrete time.

\paragraph{Transitions.} Fix $q\in(0,1)$:
\begin{equation}
P_{i,i+1}=q,
\qquad
P_{i,i-1}=1-q .
\end{equation}

\paragraph{Moments and law.} $\ex{X_n}=X_0+n(2q-1)$,
$\var{X_n}=4nq(1-q)$; for deterministic $X_0$,
$\pr{X_n=X_0+2r-n}=\binom{n}{r}q^r(1-q)^{n-r}$, $r=0,1,\ldots,n$.

\paragraph{Anchors.} The embedded jump chain of M5 has
$q=\lambda/(\lambda+\mu)$. With absorbing barriers $0$ and $M$, the hitting
probability of $M$ before $0$ is
$h_i=(1-\theta^i)/(1-\theta^M)$ with $\theta=(1-q)/q$, reducing to $i/M$ at
$q=\tfrac12$.

\section{M2: Galton--Watson branching process (general and binary)}
\source{\S2.2.1.2; developed throughout thesis Chapter 3}

\paragraph{State space.} $\{0,1,2,\ldots\}$, discrete time.

\paragraph{Transitions.} With offspring law $\pr{L=k}=p_k$ and mean
$m=\ex{L}$, from $Z_0=1$:
\begin{equation}
Z_{n+1}=\sum_{i=1}^{Z_n}L_i,
\end{equation}
the $L_i$ independent copies of $L$. Under the probability generating
function $\phi(z)=\ex{z^L}$ the generational PGFs compose:
$\phi_{n+1}=\phi_n\circ\phi$, $\phi_0(z)=z$.

\paragraph{Binary case used throughout the thesis.}
\begin{equation}
p_2=p,\qquad
p_0=1-p,\qquad
p_k=0\ \ (k\notin\{0,2\}),
\end{equation}
with $m=2p$, $\ex{Z_n}=(2p)^n$, $\phi(z)=pz^2+(1-p)$, and the survival
recursion
\begin{equation}
S_{n+1}=2pS_n-pS_n^2,
\qquad S_0=1,
\end{equation}
where $S_n=\pr{Z_n>0}$. For $0<p\le\tfrac12$ survival decays to zero;
for $p>\tfrac12$ it decreases to the fixed point $(2p-1)/p$.

\paragraph{Chapter 3 apparatus.} For $0<p<\tfrac12$ the constant
$A(p)$ is defined as the limit of $S_n/(2p)^n$ (existence via a convergent
infinite product; degenerate endpoint $A(0)=\tfrac12$); the population
conditioned on survival settles to a quasi-stationary law with mean
$1/A(p)$; near criticality $A(p)\sim 2(1-2p)$; rescaling the survival
recursion to the logistic map identifies $A(p)$ with a value of the
associated Koenigs linearising function.

\section{M3: Poisson process}
\source{\S2.2.2.2}

\paragraph{State space.} $\{0,1,2,\ldots\}$; sole non-zero transition rate
\begin{equation}
q_{n,n+1}=\alpha,
\qquad n\ge0 .
\end{equation}

\paragraph{Law.} From $N_0=0$:
$\pr{N_t=n}=\ee^{-\alpha t}(\alpha t)^n/n!$, with
$\ex{N_t}=\var{N_t}=\alpha t$. Thinning with retention probability
$\varpi$ gives a Poisson process of rate $\varpi\alpha$.

\section{M4: Yule process}
\source{\S2.2.2.3}

\paragraph{State space.} $\{1,2,\ldots\}$ from a single founder; sole
non-zero transition rate
\begin{equation}
q_{n,n+1}=\lambda n,
\qquad n\ge1 .
\end{equation}

\paragraph{Law.} $\pr{Y_t=n}=\ee^{-\lambda t}\lrb{1-\ee^{-\lambda t}}^{\,n-1}$,
$\ex{Y_t}=\ee^{\lambda t}$,
$\var{Y_t}=\ee^{\lambda t}\lrb{\ee^{\lambda t}-1}$; from $k$ founders, the
sum of $k$ independent one-founder processes (negative binomial).

\section{M5: Linear birth--death process, and time-inhomogeneous rates}
\source{\S2.2.2.4; \S2.2.3 paragraph ``Linear birth--death rates''}

\paragraph{State space.} $\{0,1,2,\ldots\}$; state $0$ absorbing. Non-zero
rates out of $n\ge1$:
\begin{equation}
q_{n,n+1}=\lambda n,
\qquad
q_{n,n-1}=\mu n .
\end{equation}
$\mu=0$ recovers the Yule process (M4).

\paragraph{Anchors.} From $X_0=N$ the mean is $N\ee^{(\lambda-\mu)t}$.
Extinction is certain for one founder when $\lambda\le\mu$ and has
probability $\mu/\lambda$ when $\lambda>\mu$; from $N$ founders,
$(\mu/\lambda)^N$. The embedded chain is M1 with $q=\lambda/(\lambda+\mu)$.

\paragraph{Time-inhomogeneous variant.} Per-capita rates $\lambda(t)$,
$\mu(t)$. The mean closes:
$\ex{X_t\mid X_0=n_0}=n_0\exp\lrb{\int_0^t(\lambda(r)-\mu(r))\,\dd r}$.
The backward extinction probability
$u(s,T)=\pr{X_T=0\mid X_s=1}$ satisfies the backward Riccati equation
\begin{equation}
-\pdd{u}{s}
=\mu(s)-\bigl(\lambda(s)+\mu(s)\bigr)u+\lambda(s)u^2,
\qquad
u(T,T)=0,
\end{equation}
whereas the forward-endpoint probability is
$u_{\mathrm f}(t)=B(t)/(1+B(t))$ with
$B(t)=\int_0^t\mu(s)\exp\lrb{\int_0^s(\mu(r)-\lambda(r))\,\dd r}\,\dd s$;
the two endpoints must not be conflated.

\section{M6: The birth--death--catastrophe (BDC) process}
\source{Definition in \S2.2.2.5; restated in \S4.2; canonical restatement
and notation in \S5.2.1--5.2.2. Conventions: \S4.2.3--4.2.4. Variant:
several founders, \S5.9.}

\paragraph{State space.} $\{0,1,2,\ldots\}\cup\{H\}$: the load within a
compartment while it remains intact, with $H$ a distinct absorbing rupture
state. Both $0$ and $H$ are absorbing.

\paragraph{Transitions.} Each unit independently gives birth at rate
$\lambda$, dies at rate $\mu$, and triggers catastrophe (rupture) at rate
$\delta$. For $n\ge1$ the non-zero transition rates are
\begin{equation}
q_{n,n+1}=\lambda n,
\qquad
q_{n,n-1}=\mu n,
\qquad
q_{n,H}=\delta n,
\end{equation}
and, in first-step form,
\begin{align}
\pr{\Delta \Xt = 1,\ \Delta \Wt = 0 \mid \Xt=n} &= \lambda n\,\Delta t+o(\Delta t)
&&\text{(birth)}\notag\\
\pr{\Delta \Xt = -1,\ \Delta \Wt = 0 \mid \Xt=n} &= \mu n\,\Delta t+o(\Delta t)
&&\text{(death)}\\
\pr{X_{t+\Delta t} = H,\ W_{t+\Delta t} = n \mid \Xt=n} &= \delta n\,\Delta t+o(\Delta t)
&&\text{(catastrophe)}\notag\\
\pr{\Delta \Xt = 0,\ \Delta \Wt = 0 \mid \Xt=n} &=
1-(\lambda+\mu+\delta)n\,\Delta t+o(\Delta t)
&&\text{(no event).}\notag
\end{align}
Default initial condition: $\Xt=1$ (working point
$(\lambda,\mu,\delta)=(1,0.2,0.05)$ in the Chapter 4 figures).

\paragraph{Release counter.} With
$\tau=\inf\{t:\Xt=H\}$ the catastrophe time,
\begin{equation}
\Wt=
\begin{cases}
0,&t<\tau,\\
X_{\tau-},&t\ge\tau,
\end{cases}
\qquad
\Kr:=X_{\tau-}
\end{equation}
is the released count and $\Kr$ the release size at rupture. The left limit
is essential.

\paragraph{Two conventions for rupture (\S4.2.3).} Sending rupture to $0$
instead of $H$ is \emph{killing}; sending it to $H$ is \emph{catastrophe}
(the thesis's convention). Killing makes internal extinction and rupture
indistinguishable; with $H$ separate one may distinguish the
internal-extinction probability $D$, the rupture probability $1-I$, and the
productive probability $\Ifix=I-D$.

\paragraph{Numerical-load convention (\S4.2.4).} Wherever $\Xt$ appears
arithmetically, $H\mapsto 0$; $\Kr$ (release size) is unrelated to
$K(t)=\ex{\Xt^2}$ (second moment).

\paragraph{Characteristic roots and fixation functions (\S5.2).} With
$\eta=(\lambda+\mu+\delta)/(2\lambda)$,
\begin{equation}
a=\eta+\sqrt{\eta^2-\tfrac{\mu}{\lambda}},
\qquad
b=\eta-\sqrt{\eta^2-\tfrac{\mu}{\lambda}},
\qquad
b<1<a,
\end{equation}
$A:=a-1$, $B:=1-b$, $\theta:=\lambda(a-b)
=\sqrt{(\lambda+\mu+\delta)^2-4\lambda\mu}$, and $AB=\delta/\lambda$.
The no-catastrophe and internal-extinction probabilities
\begin{equation}
I(t)=\pr{\text{no catastrophe by }t},
\qquad
D(t)=\pr{\text{internal extinction by }t,\ \text{no catastrophe}},
\qquad
\Ifix(t)=I(t)-D(t)
\end{equation}
both satisfy the Riccati equation
$\ddt{I}=\lambda(I-a)(I-b)=\mu-(\lambda+\mu+\delta)I+\lambda I^2$,
differently initialised ($I(0)=1$, $D(0)=0$), with closed forms
$I=\frac{aB+bAw}{B+Aw}$, $D=\frac{ab(w-1)}{aw-b}$,
$\Ifix=\frac{(a-b)^2w}{(B+Aw)(aw-b)}$ at $w=\ee^{\theta t}$.
$b$ is the probability of clearing internally without ever rupturing.
Anchors used later: $V_\infty=\ex{\Wt(\infty)}=\frac{a(1-b)}{a-1}=\frac{aB}{A}$,
$\ex{\Kr\mid\text{catastrophe}}=a/(a-1)$, and the mean productive lifetime
$\mean{T_{\rm prod}}=\lambda^{-1}\log\frac{a}{a-1}$.

\paragraph{Several founders (\S5.9).} From founding multiplicity $k$, the
fixation functions become $I(\alpha)^k-D(\alpha)^k$ with release kernel
$g_k(\alpha)=\delta K_k(\alpha)$; families founded independently factor
until the first catastrophe, which stops all lineages at once.

\paragraph{Provenance note.} The classical treatments cited in \S5.2.2
allow immigration and distributed catastrophe sizes; the thesis specialises
to the constant-rate, unit-catastrophe case above. A state-dependent
catastrophe is not a mean-field hazard:
$\pr{\text{no catastrophe before generation }n}
=\ex{(1-\chi)^{\mathcal E_n}}$ over cumulative exposure
$\mathcal E_n=\sum_{i<n}\widetilde X_i$, and Jensen's inequality forbids
replacing $\mathcal E_n$ by its mean (\S2.2.2.5, Remark).

\section{M7: Seasonal two-population process}
\source{\S2.2.3.1}

\paragraph{State space.} $\{0,1,\ldots\}^2$: counts $X_t$, $Y_t$ of two
interacting populations.

\paragraph{Transitions.} Time-inhomogeneous generator
\begin{equation}
q_{(x,y),(x+1,y)}(t)=b\mathcal S^+(t)\,x,
\qquad
q_{(x,y),(x-1,y+1)}(t)=exy,
\qquad
q_{(x,y),(x,y-1)}(t)=dy,
\end{equation}
with seasonal multiplier
$\mathcal S^+(t)=\max\{0,\,1+A\sin(\omega t)\}$ (clipped because a rate
cannot be negative). For $A>1$ each period $2\pi/\omega$ contains a closed
season of length $(2/\omega)\arccos(1/A)$. On the axis $y=0$ the $X$
population is a time-inhomogeneous Yule process; in the interior the first
moments involve $\ex{X_tY_t}$ and do not close. The deterministic
comparison system is $\dot x=b\mathcal S^+(t)x-exy$, $\dot y=exy-dy$.
Exact paths are simulated by thinning against $\mathcal S^+(t)\le 1+A$.

\section{M8: Logistic-environment type-count process}
\source{\S2.2.3.2; full distribution in \S9.2}

\paragraph{Driving environment.} Total abundance $N(t)$ follows the
logistic ODE
$\dda{N}{t}=\varrho N\lrb{1-\frac{N}{K}}$ (Chapter 9 notation; Chapter 2
uses $r$, $C$), with solution
$N(t)=K\ee^{\varrho t}/(\xi+\ee^{\varrho t})$, $\xi=K/N_0-1$.

\paragraph{State space and transitions.} $S_t\in\{1,2,\ldots\}$, $S_0=1$, no
removal; a pure birth process with time-dependent per-capita rate
\begin{equation}
q_{s,s+1}(t)=\beta(t)\,s,
\qquad
\beta(t)=\sigma\bigl(K-N(t)\bigr)
=\sigma K\,\frac{\xi}{\xi+\ee^{\varrho t}} .
\end{equation}
Exactly a Yule process under a deterministic time change; one-way coupling
(the type count does not feed back on $N$).

\paragraph{Law.} With $\Bint(t)=\int_0^t\beta$, the marginal law is
geometric, $\pr{S_t=n}=\ee^{-\Bint(t)}(1-\ee^{-\Bint(t)})^{n-1}$ for
$n\ge1$ (\S9.2.4); $\ex{S_t}=\ee^{\Bint(t)}$ converges to the finite limit
$(K/N_0)^{\sigma K/\varrho}$ (Chapter 2: $(C/N_0)^{\sigma C/r}$), and
$\Bint(\infty)=\frac{\sigma K}{\varrho}\ln\frac{K}{N_0}$, so $S_\infty$ is
a proper geometric limit.

\paragraph{Related numerical variant (\S2.2.3.2).} A
logistic-environment birth--death comparison with origination
$\beta_{\mathrm{bd}}(t)=\sigma N'(t)$ and constant per-type removal
$\varepsilon$ is used for figure-level comparison only.

\section{M9: Accumulation with gated release (coupled ODE--CTMC, two-way)}
\source{\S2.2.4.1}

\paragraph{State space.} $(A_t,R_t,G_t)\in[0,\infty)\times[0,\infty)\times
\{0,1\}$: accumulated amount, transferred amount, gate. A
piecewise-deterministic Markov process.

\paragraph{Flow (between gate jumps).}
\begin{equation}
\dda{A_t}{t}=\rho-G_t\,cA_t,
\qquad
\dda{R_t}{t}=G_t\,cA_t ,
\end{equation}
with external input $\rho>0$ and transfer rate $c>0$.

\paragraph{Jumps.} The gate opens at a state-dependent rate and closes at a
constant one:
\begin{equation}
\pr{G_{t+\Delta t}=1\mid G_t=0,A_t=a}
=\eta a\,\Delta t+o(\Delta t),
\qquad
\pr{G_{t+\Delta t}=0\mid G_t=1}
=\gamma\,\Delta t+o(\Delta t).
\end{equation}
Pathwise balance: $A_t+R_t=A_0+R_0+\rho t$. The moment
$\dda{}{t}\ex{A_t}=\rho-c\ex{G_tA_t}$ is not closed.
Illustration parameters: $\rho=2$, $c=4$, $\eta=0.015$, $\gamma=1.45$,
$(A_0,R_0,G_0)=(5,0,0)$.

\section{M10: Accumulation with surging gated release (mode-driven)}
\source{\S2.2.4.2}

\paragraph{Specification.} The flow and closing rate of M9 are retained;
opening is autonomous,
\begin{equation}
\pr{G_{t+\Delta t}=1\mid G_t=0}
=\nu\,\Delta t+o(\Delta t),
\end{equation}
so $G_t$ is an autonomous two-state CTMC. The gate mean closes:
$\ex{G_t}\to\nu/(\nu+\gamma)$; the balance law is unchanged.
Illustration parameters: as M9 with $\nu=0.3$.

\medskip
\noindent\textit{Context.} \S2.2.4 also states the three-way taxonomy of
couplings (chain switching the vector field; prescribed trajectory driving
the rates; two-way) with joint generator
$(\mathcal Lf)(\mathbf y,i)=F(\mathbf y,i)\cdot\nabla_{\mathbf y}f
+\sum_{j\ne i}q_{ij}(\mathbf y)[f(\mathbf y,j)-f(\mathbf y,i)]$; M9 is the
two-way case and M10 the first case.

\section{M11: Absorption-only model}
\source{Appendix 2.B, \S2.B.1}

\paragraph{Setting.} One compartment drawing particles from a surrounding
medium; $\Xt$ particles inside, $Y_t$ outside; $X_0=0$, $Y_0=N$.

\paragraph{Transitions.}
\begin{equation}
\pr{\Delta X=1,\ \Delta Y=-1\mid Y_t=j}=j\alpha\,\Delta t
\qquad\text{(absorption)}
\end{equation}
and no event otherwise. Particles are neither created nor destroyed
($p_{i,j}(t)=0$ unless $i+j=N$).

\paragraph{Law.} $X_t\sim\operatorname{Binomial}(N,\,1-\ee^{-\alpha t})$;
single-particle PGF
$G_1=(1-\ee^{-\alpha t})x+\ee^{-\alpha t}y$.

\section{M12: Absorption--death model}
\source{\S2.3.3.2 worked example}

\paragraph{Transitions.} To the absorption channel of M11 add interior
death:
\begin{align*}
\pr{\Delta X=1,\ \Delta Y=-1\mid Y_t=j} &= j\alpha\,\Delta t, &&\text{(absorption)}\\
\pr{\Delta X=-1,\ \Delta Y=0\mid X_t=i} &= i\mu\,\Delta t, &&\text{(death)}
\end{align*}
with $X_0=0$, $Y_0=N$. Single-particle state probabilities:
$p_{0,1}(t)=\ee^{-\alpha t}$,
$p_{1,0}(t)=\frac{\alpha}{\mu-\alpha}(\ee^{-\alpha t}-\ee^{-\mu t})$,
$p_{0,0}=1-p_{0,1}-p_{1,0}$; the $N$-particle PGF is the $N$th power of the
single-particle one. Generator-level PDE:
$\pddt{G}=\mu(1-x)\pdd{G}{x}+\alpha(x-y)\pdd{G}{y}$.

\section{M13: Absorption--birth--death model}
\source{Appendix 2.B, \S2.B.2--2.B.3; closed form in \S2.B.2 Proposition}

\paragraph{Transitions.} Absorption as M11/M12 plus interior birth:
\begin{align*}
\pr{\Delta X=1,\ \Delta Y=-1\mid Y_t=j} &= j\alpha\,\Delta t, &&\text{(absorption)}\\
\pr{\Delta X=1,\ \Delta Y=0\mid X_t=i} &= i\lambda\,\Delta t, &&\text{(birth)}\\
\pr{\Delta X=-1,\ \Delta Y=0\mid X_t=i} &= i\mu\,\Delta t, &&\text{(death)}
\end{align*}
with $X_0=0$, $Y_0=N$.

\paragraph{Closed form.} With $x_+=1$, $x_-=\mu/\lambda$,
$b_1=\lambda-\mu$, $\omega=\alpha/b_1$, $W=(x-x_+)/(x-x_-)$ and
$\Psi(\zeta)=x_++\frac{\alpha}{\lambda}\sum_{n\ge1}\frac{\zeta^n}{n+\omega}
=x_++(x_+-x_-)\lrb{{}_2F_1(1,\omega;\omega+1;\zeta)-1}$,
\begin{equation}
G(x,y,t)=\Bigl\{\ee^{-\alpha t}\bigl[y-\Psi(W)\bigr]
+\Psi\bigl(W\ee^{b_1t}\bigr)\Bigr\}^{N}.
\end{equation}
The parameter-regular convolution representation is
$g(x,y,t)=\ee^{-\alpha t}y+\alpha\ee^{-\alpha t}\int_0^t\ee^{\alpha
v}F(x,v)\,\dd v$ with $F$ the interior birth--death PGF, and
$G=g^N$. Resonance at $\alpha=n(\mu-\lambda)$, $n=1,2,\ldots$; the
generating function stays regular there.

\section{M14: Chained immediate transfer ($\mu=0$ birth--catastrophe chain)}
\source{\S5.10; described as ``successive birth--death--catastrophe'' in
\S8.2.1}

\paragraph{Specification.} $M$ compartments, all empty but one holding a
single particle. Each runs the $\mu=0$ BDC process (M6 with $\mu=0$, so
only births $\lambda n$ and catastrophe $\delta n$); on rupture the entire
released load passes immediately and wholesale into a fresh compartment,
founding the next process with that load as initial census. No extracellular
phase, no clearance, no re-founding; all $M$ ruptures occur almost surely.

\paragraph{Anchors.} With $s=\delta/(\lambda+\delta)$ and
$\rho=\lambda/(\lambda+\delta)=1-s$: event types ignore the load
(birth with probability $\rho$, catastrophe with probability $s$), the
births before catastrophe in a compartment founded with any load are
$G\sim\operatorname{Geom}_0(s)$, and the $k$th release size is
$r_k=1+G_1+\cdots+G_k$ with
\begin{equation}
\pr{r_k=n}=\binom{n+k-2}{k-1}s^k\rho^{\,n-1},
\qquad n\ge1 .
\end{equation}
The case $\mu>0$ is recorded as open.

\section{M15: Budding comparator compartment}
\source{\S5.11; restated \S6.2.2}

\paragraph{Specification.} A budding compartment experiences two
independent Poisson event streams: release of one particle at rate $p$ and
removal of the compartment at rate $d_{\Act}$, continuing until removal.

\paragraph{Anchors.}
\begin{equation}
S_{\rm bud}(\alpha)=\ee^{-d_{\Act}\alpha},
\qquad
V_{\rm bud}(\alpha)=\frac{p}{d_{\Act}}\lrb{1-\ee^{-d_{\Act}\alpha}},
\end{equation}
and the total number released is geometric on $\{0,1,2,\ldots\}$:
\begin{equation}
\pr{\Kr_{\rm bud}=k}=\frac{d_{\Act}}{d_{\Act}+p}
\lrb{\frac{p}{d_{\Act}+p}}^{k},
\qquad k=0,1,2,\ldots .
\end{equation}
Structural contrast with rupture-based release (M6): here release and
removal are independent events; in M6 they are the same event.

% ----------------------------------------------------------------------------
\chapter{The population layer (thesis Chapter 6)}
% ----------------------------------------------------------------------------

\section{M16: Standard two-variable population model (deterministic projection)}
\source{\S6.2.1; three-variable version in \S1.5}

\paragraph{Specification (fixed target count).} Active compartments $\Act$
and free particles $\Free$ obey
\begin{equation}
\ddt{\Act}=\gamma T\,\Free-d_{\Act}\,\Act,
\qquad
\ddt{\Free}=p\,\Act-c\,\Free,
\qquad
R_0=\frac{\gamma T\,p}{c\,d_{\Act}},
\end{equation}
with target sites $T$ fixed, founding rate $\gamma$, particle clearance $c$,
release rate $p$, removal rate $d_{\Act}$. \emph{Deterministic}; the thesis
shows it is the exact one-compartment limit of M15 and a projection of the
renewal system M17, not a Markov model itself. \S1.5 gives the full
three-variable version
$\ddt{T}=s-\gamma TV$, $\ddt{I}=\gamma TV-d_I I$, $\ddt{V}=pI-cV$.

\section{M17: Kernel-driven renewal system}
\source{Kernels \S6.3; system and exactness \S6.4.1--6.4.2; skeleton \S6.4.3}

\paragraph{Intracompartment layer.} The load inside a compartment of age
$\alpha$ follows M6. Two kernels:
\begin{equation}
S(\alpha):=\Ifix(\alpha)
=\frac{(a-b)^2\,\ee^{\theta\alpha}}
{\lrb{B+A\ee^{\theta\alpha}}\lrb{a\ee^{\theta\alpha}-b}},
\qquad
g(\alpha):=\delta K(\alpha)=V'(\alpha),
\end{equation}
the survival kernel (probability still productive at age $\alpha$) and the
release kernel (expected release flux at age $\alpha$);
$g$ is explicit because a rupture fires at rate $\delta X_\alpha$ and
releases $X_\alpha$, giving $\delta\mean{X_\alpha^2}\dd\alpha$.

\paragraph{Population system.} With incidence $i(t)=\gamma T\Free(t)$ and
$\Act_0$ initially active compartments of age zero:
\begin{align}
\Act(t)&=\Act_0\,\Ifix(t)+\int_0^t i(t-\alpha)\,\Ifix(\alpha)\,\dd\alpha
=\Act_0\,\Ifix+(i*\Ifix)(t),\\
\ddt{\Free}(t)
&=\Act_0\,g(t)+\int_0^t i(t-\alpha)\,g(\alpha)\,\dd\alpha-c\,\Free(t)
=\Act_0\,g+(i*g)(t)-c\,\Free(t).
\end{align}
These are the first moments \emph{exactly} of the stochastic process in
which each founding event starts an independent BDC compartment and each
particle founds at rate $\gamma T$ or clears at rate $c$ (linear incidence
is what closes the hierarchy).

\paragraph{Skeleton.} Any kernel pair $(S,g)$ induces
\begin{equation}
\Act(t)=\Act_0\,S(t)+(i*S)(t),
\qquad
\ddt{\Free}=\Act_0\,g(t)+(i*g)(t)-c\,\Free(t),
\end{equation}
with effective parameters $p_{\rm eff}(r)=\Lap{g}(r)/\Lap{S}(r)$ and, for
linear incidence, $R_0=(\gamma T/c)\int_0^\infty g$. Standing assumptions
(A1)--(A5): constant $T$; continuum convention (particles not consumed on
founding); synchronous zero-age founding cohort at $t=0$; independent
compartments, single founding; no spatial structure.

\section{M18: Branching approximation of the early phase; the flooding parameter}
\source{\S6.6.1--6.6.2}

\paragraph{Specification.} Each released particle independently founds a
new compartment at rate $\gamma T$ or is cleared at rate $c$, hence founds
with probability
\begin{equation}
q=\frac{\gamma T}{\gamma T+c}.
\end{equation}
The offspring of one active compartment is the release size $\Kr$ (M6)
thinned by $q$: offspring PGF
\begin{equation}
G_{\rm off}(z)=b+\frac{\delta}{\lambda}\,\frac{y}{a-y},
\qquad y=1-q+qz,
\end{equation}
including the mass $b$ of compartments that never rupture. Mean
$m=qV_\infty$ ($\approx R_0$ when $\gamma T\ll c$). The comparator (M15)
is matched at equal total output $p/d_{\Act}=V_\infty$; its thinned
offspring law is geometric with the same mean.

\paragraph{Flooding parameter and criterion.}
\begin{equation}
L:=a(1-b)=a-\frac{\mu}{\lambda},
\qquad\text{so that}\qquad
V_\infty=\frac{L}{a-1}.
\end{equation}
For $m>1$ the establishment extinction probabilities satisfy
$z_{\rm ext}^{\rm rupture}-z_{\rm ext}^{\rm bud}=(L-1)\lrb{\frac1m-1}$, so
one-shot release is advantageous if and only if
\begin{equation}
\delta(\lambda+\mu)>\lambda\mu
\qquad\Longleftrightarrow\qquad
\frac{1}{\delta}<\frac{1}{\lambda}+\frac{1}{\mu}
\end{equation}
(with the condition holding for all $\delta>0$ when $\mu=0$).

\section{M19: Reset process (recurrent release)}
\source{\S6.7.1}

\paragraph{Specification.} While the compartment is alive, the load
undergoes birth at rate $\lambda n$, death at rate $\mu n$, optional
immigration at rate $\nu$, and a \emph{dump that resets}: at rate
$\delta n$ the whole pool is released and the load returns to zero, after
which it may re-accumulate; the compartment is removed at a separate,
load-independent rate $d_{\Act}$. Dumps do not absorb.

\paragraph{Kernels and anchors.}
\begin{equation}
S(\alpha)=\ee^{-d_{\Act}\alpha},
\qquad
g(\alpha)=\delta\,\ee^{-d_{\Act}\alpha}\,\mean{X_\alpha^2},
\qquad
p_{\rm eff}(r)=\lrb{r+d_{\Act}}\Lap{g}(r),
\qquad
d_{\Act,{\rm eff}}\equiv d_{\Act},
\end{equation}
with $\mean{X_\alpha^2}$ computed for the non-absorbing reset process.
$R_0=(\gamma T/c)\int_0^\infty g$ is finite only for $d_{\Act}>0$; in
the mature-compartment limit, $p_\infty=\delta\,\mathbb{E}_\pi[X^2]$ for
stationary law $\pi$.

\section{M20: Partial release (boolean thinning)}
\source{\S6.7.2}

\paragraph{Specification.} Release events fire on the load-dependent clock
$\delta X$ and remove each unit independently with probability $\varphi$
(the default); at mean-field level, with $Q$ the total intracompartmental
stock across productive compartments,
\begin{equation}
\ddt{\Act}=\gamma T\,\Free-d_{\Act}\Act,
\qquad
\ddt{Q}=\nu\Act-\delta\varphi\,\frac{Q^2}{\Act}-d_{\Act}Q,
\qquad
\ddt{\Free}=\delta\varphi\,\frac{Q^2}{\Act}-c\,\Free.
\end{equation}
With a load-independent clock $\lambda(X)=\delta$ instead, mean export is
$\delta\varphi Q$ and the equations collapse to stock-and-export dynamics
with $\varepsilon_{\rm eff}=\delta\varphi$: partial release is invisible in
the mean field exactly when the clock is load-independent; in single-compartment
release-count distributions it is always visible.

\section{M21: Budding-kernel extension with sigmoid incidence}
\source{\S6.7.3}

\paragraph{Specification.} A non-catastrophe application: the mean load
follows an immigration-like $\bar x'\approx\nu-\mu\bar x$. The skeleton of
M17 carries the application on a different kernel pair $(S_{\rm b},g_{\rm
b})$ --- probability of still being in a productive lineage, and density of
productive sojourn times --- with a multi-stage delay (near-zero release for
small $\alpha$, then rising and falling) supplying the eclipse; the
effective-parameter map carries over verbatim. Near the observed
establishment threshold the data favour a sigmoid incidence:
\begin{equation}
f(\Free,T)=\gamma T\,\Free\cdot\frac{\Free}{K_A+\Free},
\end{equation}
under which the first-moment exactness of M17 no longer holds.

% ----------------------------------------------------------------------------
\chapter{Two-type extension (thesis Chapter 7)}
% ----------------------------------------------------------------------------

\section{M22: Two-type birth--death--conversion with type-specific catastrophe rates}
\source{\S7.2.1 (definition and rate list), \S7.2.2 (occupation-time
representation), \S7.2.4 (dimensionless form)}

\paragraph{State space.} $(x,y)\in\mathbb Z_{\geq0}^2$ before absorption;
a single global absorbing state $(\R,\R)$.

\paragraph{Transitions.} Type 1 gives birth at rate $\lambda_1$, dies at
rate $\mu_1$, and converts irreversibly to type 2 at rate $\nu$; type 2
gives birth at rate $\lambda_2$ and dies at rate $\mu_2$; all rates
non-negative. Before absorption, the non-zero transition rates from
$(x,y)$ are
\begin{equation}
\begin{array}{@{}r@{\;}l@{\qquad}l@{}}
(x,y)&\longrightarrow(x+1,y)&\text{at rate}\quad\lambda_1x,\\
(x,y)&\longrightarrow(x-1,y)&\text{at rate}\quad\mu_1x,\\
(x,y)&\longrightarrow(x-1,y+1)&\text{at rate}\quad\nu x,\\
(x,y)&\longrightarrow(x,y+1)&\text{at rate}\quad\lambda_2y,\\
(x,y)&\longrightarrow(x,y-1)&\text{at rate}\quad\mu_2y,\\
(x,y)&\longrightarrow(\R,\R)&\text{at rate}\quad\delta_1x+\delta_2y,
\end{array}
\end{equation}
the last being the catastrophe: one unit contributes a type-specific rate,
and the first trigger absorbs the whole process. One reading assigns the
types to two internal states of a unit within a compartment, catastrophe
representing irreversible loss of containment.

\paragraph{Object of study.}
$S(t)=\Prob_{(1,0)}\{\tau_c>t\}$ and
$G(t)=\Prob_{(0,1)}\{\tau_c>t\}$: no-catastrophe probabilities, not
population-survival probabilities. They admit the occupation-time
representation
\begin{equation}
S(t)=\E_{(1,0)}\!\left[
\exp\!\left\{-\int_0^t(\delta_1X_s+\delta_2Y_s)\,\dd s\right\}\right],
\qquad
G(t)=\E_{(0,1)}\!\left[
\exp\!\left\{-\int_0^t(\delta_1X_s+\delta_2Y_s)\,\dd s\right\}\right],
\end{equation}
and satisfy the triangular backward system
\begin{align}
\frac{\dd S}{\dd t}
&=\lambda_1S^2-(\lambda_1+\mu_1+\nu+\delta_1)S+\mu_1+\nu G,
&S(0)&=1,\ S'(0)=-\delta_1,\\
\frac{\dd G}{\dd t}
&=\lambda_2G^2-(\lambda_2+\mu_2+\delta_2)G+\mu_2,
&G(0)&=1,\ G'(0)=-\delta_2.
\end{align}
Dimensionless form: time in units of $\lambda_1$; rates
$\hat\mu_i=\mu_i/\lambda_1$, $\hat\nu=\nu/\lambda_1$,
$\hat\delta_i=\delta_i/\lambda_1$, $\hat\lambda_2=\lambda_2/\lambda_1$.
The main formula assumes $\hat\lambda_2>0$, $\hat\delta_2>0$,
$\hat\nu>0$; the zero-birth and zero-catastrophe boundaries are treated
separately, and with the appropriate initial data the zero-catastrophe case
recovers the two-type branching equations of Antal and Krapivsky.
\S7.5 compares four catastrophe-rate regimes (per-type constant, type-2
only, and mixtures) exactly.

% ----------------------------------------------------------------------------
\chapter{Depletion models (thesis Chapter 8)}
% ----------------------------------------------------------------------------

\section{M23: Replicator--suppressor process}
\source{\S8.3.2 (definition); first-step analysis \S8.3.3; fates \S8.3.4}

\paragraph{Setting.} A compartment contains a replicating population and a
finite store of suppressor units; suppressors are consumed one-for-one as
they act, and the store is never replenished.

\paragraph{State space.}
$\mathcal{S}=\{(i,j): i\in\mathbb{N}\cup\{0\},\ j\in\{0,\dots,N\}\}$:
$\Xt$ replicating units in the compartment, $Q_t$ remaining suppressor
units (store size $N$ fixed at founding).

\paragraph{Transitions.} For $(i,j)\in\mathcal S$:
\begin{align}
\pr{\Delta X_t=1,\;\Delta Q_t=0 \mid X_t=i,\,Q_t=j}
  &= \lambda i\,\Delta t+o(\Delta t)
  &&\text{(birth)},\\
\pr{\Delta X_t=-1,\;\Delta Q_t=-1 \mid X_t=i,\,Q_t=j}
  &= \varsigma ij\,\Delta t+o(\Delta t)
  &&\text{(suppression)},\\
\pr{\Delta X_t=0,\;\Delta Q_t=0 \mid X_t=i,\,Q_t=j}
  &= 1-\lrb{\lambda i+\varsigma ij}\Delta t+o(\Delta t)
  &&\text{(no event)}.
\end{align}
The product $\varsigma ij$ is the only nonlinearity. The column $\Xt=0$ is
extinction; the row $Q_t=0$ is a spent store, from which the process is a
pure birth process and survival is certain. Interest centres on
$\pr{\text{survival}}=1-\lim_{t\to\infty}\pr{X_t=0}$; survival is certain
whenever the founding cohort outnumbers the store ($i>j$), and in the
contested region $1\le i\le j$ the survival map $\pi_{i,j}$ is computed by
first-step analysis.

\section{M24: Fixed-removal extension of the chained process}
\source{\S8.4.1--8.4.2; assumptions in \S8.4.1 Remarks}

\paragraph{Specification.} The chained-rupture picture (M14's complete
transfer, or dispersal to new compartments) is extended by one assumption:
after each release event, a fixed count of units is removed. The removal
rule is
\begin{equation}
\text{removal} =
\begin{cases}
r, & r\le\vartheta,\\
\vartheta, & r>\vartheta,
\end{cases}
\qquad\text{that is, } \min(r,\vartheta),
\end{equation}
with $r$ the release size and $\vartheta$ the fixed removal budget
(an idealisation of a saturating casualty law
$\vartheta(r)\sim\log(r+1)$). Under the framing used for the calculation,
each surviving released unit founds a new compartment.

\paragraph{Assumptions.} Remover-agent prevalence resets between release
events (steady state, uniform distribution --- this makes $\vartheta$ the
same at every delivery and the independence exact); the release-size
distribution is that of M6 including internal-death outcomes
($\ex{r}=V_\infty$).

\paragraph{Quantity of interest.} The expected replication ratio
\begin{equation}
\mathcal{R}(\vartheta)=\ex{\max(0,\,r-\vartheta)}
=\bigl[\ex{r\mid r\ge\vartheta}-\vartheta\bigr]\pr{r\ge\vartheta},
\end{equation}
the expected number of newly founded compartments created by one founding
event. Deliberately not the $R_0$ of M17--M18; the relation is stated in
\S8.5.1.

% ----------------------------------------------------------------------------
\chapter{Models with non-constant rates (thesis Chapters 9--10)}
% ----------------------------------------------------------------------------

\section{M25: Competitive-reversal chain (evolving resistance)}
\source{\S9.3.2; critical point \S9.3.4}

\paragraph{State space.} $A_t\in\{0,1,\ldots,N\}$ in discrete time, two
populations with $A_t+B_t=N$; $i=0$ and $i=N$ absorbing. The triple
$(A_t,R_{a,t},R_{b,t})$ is Markov; $A_t$ alone is not.

\paragraph{Transition kernel.} For $0<i<N$:
\begin{equation}
\pr{A_{t+1}=j\mid A_t=i}=
\begin{cases}
q_a=\dfrac{i}{N}(1-p_b), & j=i+1,\\[6pt]
q_b=\lrb{1-\dfrac{i}{N}}(1-p_a), & j=i-1,\\[6pt]
q_c=\dfrac{i}{N}p_b+\lrb{1-\dfrac{i}{N}}p_a, & j=i,\\[4pt]
0 & \text{otherwise.}
\end{cases}
\end{equation}
Blocking capacities
\begin{equation}
p_a=\max\left\{0,\,1-\frac{R_b}{R_a}\right\},
\qquad
p_b=\max\left\{0,\,1-\frac{R_a}{R_b}\right\},
\end{equation}
so at most one of $p_a,p_b$ is non-zero. Resistances are driven up when the
population is in a minority:
\begin{equation}
R_{k,t+1}=R_{k,t}+\varepsilon\,\Delta R_{k,t},
\qquad
\Delta R_{a,t}=\max\left\{0,\lrb{\frac{N-A_t}{A_t}}^{\!\alpha}\right\},
\quad
\Delta R_{b,t}=\max\left\{0,\lrb{\frac{A_t}{N-A_t}}^{\!\alpha}\right\},
\end{equation}
with $R_{a,0}=R_{b,0}=1$ without loss (only $\varepsilon/R_0$ matters).
The critical state $q_a=q_b$ is
$i_c=N\frac{1-p_a}{2-p_a}$ ($p_a>0$) or $N\frac{1}{2-p_b}$ ($p_b>0$), and
is unstable; it is a function of the path, not the parameters alone.

\section{M26: Fitness-coupled type-origination variant}
\source{\S9.3.6; recorded as specified but never simulated}

\paragraph{Specification.} The kernel of M25 recast with fitness in place
of resistance, a type count carried alongside and acting back on fitness:
\begin{equation}
\pr{\Delta A_t=1}=\frac{A_t}{N}\,\frac{f_a}{\bar f},
\qquad
\pr{\Delta A_t=-1}=\lrb{1-\frac{A_t}{N}}\frac{f_b}{\bar f},
\end{equation}
with $\bar f$ the mean fitness across the domain; fitness updates
\begin{equation}
f_a(t+\Delta t)=f_a(t)+\varepsilon\lrb{\frac{B_t}{A_t}}^{\!\alpha}
-\gamma S^{(a)}_t,
\qquad
f_b(t+\Delta t)=f_b(t)+\varepsilon\lrb{\frac{A_t}{B_t}}^{\!\alpha}
-\gamma S^{(b)}_t,
\end{equation}
and per-population type counts evolving as birth--death processes with
constant removal $\mu_a$ and dominance-and-fitness origination
\begin{equation}
\pr{\Delta S^{(a)}_t=1}=\beta_a S^{(a)}_t\,\Delta t,
\qquad
\beta_a=\lrb{\frac{A_t}{N}}^{\!\zeta}f_a .
\end{equation}
The penalty $-\gamma S$ is the endogenous weakening of the leader;
$\zeta$ controls how firmly dominance must be established before variation
is released.

\section{M27: Rise--run--ruin--rejuvenation (RRRR)}
\source{\S9.4.2}

\paragraph{State space.} $\Xt\in\{0,1,2,\ldots\}$ replicating agencies
with an auxiliary real potential $V(t)$; the pair $(\Xt,V(t))$ is Markov,
$\Xt$ alone is not. Described in the source as ``a
birth--death--catastrophe process with catastrophe rate
$\delta X_t/V(t)$ and an auxiliary ODE for $V$''.

\paragraph{Specification.} Potential flows in and is consumed:
\begin{equation}
\ddt{V}=\phi-c\,X_t .
\end{equation}
Given $V$, the count moves by
\begin{align}
\pr{\Delta X_t=1\mid X_t}&=\lambda X_t\,\Delta t+o(\Delta t)
&&\text{(proliferation)},\notag\\
\pr{X_{t+\Delta t}=0\mid X_t}&=\delta\,\frac{X_t}{V(t)}\,\Delta t+o(\Delta t)
&&\text{(catastrophe)},\notag\\
\pr{X_{t+\Delta t}=1\mid X_t=0}&=\omega\,\Delta t+o(\Delta t)
&&\text{(revival)},\\
\pr{X_{t+\Delta t}=X_t}&=1-\lrb{\lambda X_t
+\delta\frac{X_t}{V(t)}+\omega\ind\{X_t=0\}}\Delta t+o(\Delta t)
&&\text{(no event).}\notag
\end{align}
Catastrophe is driven by per-capita scarcity; revival keeps $0$ from
absorbing, making the process recurrent. $V$ stays positive for all time
(the hazard diverges as $V\downarrow0$). Illustration parameters:
$\phi=4$, $c=0.8$, $\lambda=0.5$, $\delta=0.25$, $\omega=0.15$.

\section{M28: Multiplicative-dissipation automaton}
\source{\S9.5 (state, competition, mutation and potential, fragmentation)}

\paragraph{State space.} An $L\times L$ lattice ($L=256$; Moore
neighbourhood, periodic boundary, torus), each cell empty or occupied.
Every occupied cell carries a type and, within it, a sub-type. For each
extant type $i$: accumulated potential $V_i$ (real), occupied cells $N_i$,
extant sub-types $X_i\le N_i$. Continuous updates (potential, mutation)
advance in chunks between discrete contests (one attempt per site per
sweep).

\paragraph{Competition.} A cell $\zeta$ carrying sub-type $u$ of type $i$
has strength
\begin{equation}
w(\zeta)=V_i\times(n_\zeta+1)^{\theta}\times\ee^{-\nu|u|/\kappa_i},
\qquad
\kappa_i=\frac{N_i}{X_i},
\end{equation}
with $n_\zeta$ matching neighbours at the contest level and $|u|$ the area
of $u$; the contest goes to $\zeta$ with probability
$w(\zeta)/(w(\zeta)+w(\zeta'))$. Between sub-types of one type the
potential cancels; between distinct types the rarity factor is omitted,
$\pr{A\text{ wins}}=\frac{V_A(n_A+1)^{\theta}}
{V_A(n_A+1)^{\theta}+V_B(n_B+1)^{\theta}}$. Contests against empty ground:
fixed strength $d$, an occupied cell takes a contested empty neighbour with
probability $1-d$, empty ground never invades.

\paragraph{Mutation and potential.} Sub-types arise within a type at
type-wide rate $R^{\text{mut}}_i=m\,V_i N_i$ (per-cell rate $mV_i$).
Each occupied cell supplies potential at rate $\varphi$ and each extant
sub-type costs $c$:
\begin{equation}
\ddt{V_i}=\varphi N_i-c\,X_i ,
\qquad\text{floored at }0,
\end{equation}
so potential can fall at all only if $c>\varphi$.

\paragraph{Fragmentation.} When a type carries too many sub-types it
fragments: every sub-type becomes an independent type at once (nothing
dies; the interaction law changes). Two triggers, tested in order:
deterministically if the potential update would take $V_i$ to zero or
below; otherwise with hazard
\begin{equation}
R^{\text{frag}}_i=\delta\,\frac{X_i}{V_i},
\end{equation}
M27's catastrophe rate with the count of individuals replaced by the count
of types. No revival term is needed: the catastrophe restarts the process.

\section{M29: Birth--conversion model, and the vitality-coupled variant}
\source{Appendix 9.A; simulated but not analysed}

\paragraph{State space.} $X_t$ unconverted individuals, $Y_t$ converted.

\paragraph{Transitions.}
\begin{align}
\pr{\Delta X_t=1,\ \Delta Y_t=0\mid X_t,Y_t}
  &=\lambda X_t\,\Delta t+o(\Delta t)
  &&\text{(birth)},\\
\pr{\Delta X_t=-1,\ \Delta Y_t=1\mid X_t,Y_t}
  &=\bigl(\delta+\chi Y_t\bigr)X_t\,\Delta t+o(\Delta t)
  &&\text{(conversion)},\\
\pr{\Delta X_t=0,\ \Delta Y_t=-1\mid X_t,Y_t}
  &=\mu Y_t\,\Delta t+o(\Delta t)
  &&\text{(clearance)}.
\end{align}
$\delta$ is the spontaneous-conversion rate (linear in $X_t$) and $\chi$
the contact-conversion coefficient (bilinear $X_tY_t$); the two mechanisms
superpose in one channel. Illustration parameters:
$(\lambda,\delta,\chi,\mu)=(1,\,0.1,\,0.02,\,10)$.

\paragraph{Vitality-coupled variant.} Carry $V(t)$ with
$\ddt{V}=\phi-c\,X_t$ and replace the middle channel by
\begin{equation}
\pr{\Delta X_t=-1,\ \Delta Y_t=1\mid X_t,Y_t}
=\lrb{\frac{\delta}{V(t)}+\chi Y_t}X_t\,\Delta t+o(\Delta t),
\end{equation}
the other channels unchanged; a population that grows large drives $V$ down
and raises its own loss rate. Not simulated.

\section{M30: Public-good accumulation process}
\source{Appendix 9.B}

\paragraph{State space.} $\Xt$ units within one compartment and $P(t)$ the
public-good concentration.

\paragraph{Specification.}
\begin{align}
\ddt{P}&=g\,X_t-c,\notag\\
\pr{\Delta X_t=1\mid X_t}&=\lambda(t)\,X_t\,\Delta t+o(\Delta t),
\qquad
\lambda(t)=\lambda_0+\alpha P(t),\\
\pr{\Delta X_t=-1\mid X_t}&=\mu\,X_t\,\Delta t+o(\Delta t),\notag
\end{align}
with production rate $g$ per unit, clearance $c$, base birth rate
$\lambda_0$, public-good advantage $\alpha$, death rate $\mu$. The regime
of interest is $\mu>\lambda_0$: subcritical alone, supercritical once
enough good has accumulated.

\paragraph{Embedded Galton--Watson reading.} In a compartment-to-compartment
sequence the process is a Galton--Watson process, supercritical exactly when
$\lim_{t\to\infty}\pr{X_t>0\mid X_0=K}>\tfrac12$; the minimum founding
cohort is
\begin{equation}
K^\ast=\inf\left\{K:\ \lim_{t\to\infty}\pr{X_t>0\mid X_0=K}>\tfrac12\right\}.
\end{equation}
Illustration parameters: $(\lambda_0,\alpha,g,c)=(0.6,0.25,0.15,0.45)$, so
$P$ falls whenever $X_t<c/g=3$.

\section{M31: Three-environment birth--death ordering}
\source{\S10.3.1}

\paragraph{Specification.} Replication in each of three environments
indexed $i\in\{1,2,3\}$ in increasing order of hostility --- treated as a
simple birth--death process (M5) with rates $(\beta_i,\mu_i)$ --- with
\begin{equation}
\mu_3>\mu_2>\mu_1
\qquad\text{and}\qquad
\beta_1=\beta_2=\beta_3=\beta ,
\end{equation}
so the net rates order as
$\beta-\mu_1>\beta-\mu_2>\beta-\mu_3$. A consistency argument, not a
derivation: the ordering is inferred from the strategy it is meant to
explain. (The surrounding discussion notes that with a depletable
suppressor store, as in M23, the $\mu_i$ are in reality declining functions
of the cohort's history.)

\end{document}
